% ============================================================================
% Fichier  : partie1/P01_C01_proprietes_securite.tex
% Titre    : La Sécurité de l'Information : Objets et Propriétés
% Auteur   : MINKA MI NGUIDJOÏ Thierry Emmanuel
% Version  : 1.0 -- novembre 2026
% ============================================================================
% Base : brouillon « Livre I » (archivé). Définitions de la loi n° 2010/012,
% art. 4 (points 24, 33, 38, 48, 58), lues dans le texte officiel ; ISO/IEC
% 27000 ; principes de Saltzer et Schroeder (1975) ; hexade de Parker (1998).
% ============================================================================

\chapter{La Sécurité de l'Information : Objets et Propriétés}
\label{chap:proprietes}
\chaptermark{Objets et propriétés}

\epigraph{%
    « Sécuriser, c'est d'abord dire ce que l'on protège,\\
    et contre quelle perte. »%
}{Principe du cours}

\niveauChapitre{Fondamental}{Aucun prérequis~; le Prologue
(Chapitre~\ref{chap:prologue}) fournit le cas.}

\begin{boxobjectifs}
\begin{itemize}
    \item Définir information, système d'information, actif, sécurité de
          l'information et cybersécurité.
    \item Énoncer les propriétés de sécurité et les rattacher aux
          définitions de la loi n°~2010/012.
    \item Exprimer les besoins de sécurité d'un actif sur une échelle.
    \item Distinguer menace, vulnérabilité, impact et risque.
    \item Appliquer les principes de conception de Saltzer et Schroeder.
\end{itemize}
\end{boxobjectifs}

\bigskip

% ============================================================================
\section*{Question directrice : qu'a-t-on perdu, au juste, à Garoua~?}
% ============================================================================

Le lundi matin, M.~ETEME dit que « la sécurité a été percée ». La phrase
ne permet aucune décision. LiZbiZ a perdu l'usage de sa facturation, elle
a peut-être perdu le secret de ses données clients, et elle ne sait plus si
ses données de stock sont exactes. Ce sont trois pertes différentes, qui
appellent trois réponses différentes. Ce chapitre donne le vocabulaire
qui permet de les nommer.

\medskip
\textbf{Plan du chapitre.}
\begin{enumerate}
    \item Information, système d'information, actif (section~\ref{sec:c01-objets}).
    \item Sécurité de l'information et cybersécurité (section~\ref{sec:c01-securite}).
    \item Les propriétés de sécurité (section~\ref{sec:c01-proprietes}).
    \item Exprimer des besoins de sécurité (section~\ref{sec:c01-besoins}).
    \item Menace, vulnérabilité, risque (section~\ref{sec:c01-risque}).
    \item Les principes de conception (section~\ref{sec:c01-principes}).
\end{enumerate}

\separateur

% ============================================================================
\section{Information, Système d'Information, Actif}
\label{sec:c01-objets}
% ============================================================================

\begin{boxdef}
\termecle{Système d'information}~: selon la \loicam{n°~2010/012}
(art.~4, point~72), « dispositif isolé ou groupe de dispositifs
interconnectés ou apparentés, assurant par lui-même ou par un ou plusieurs
de ses éléments, conformément à un programme, un traitement automatisé de
données ».

\medskip
\termecle{Actif}~: tout élément qui a de la valeur pour l'organisation.
On distingue les \emph{actifs essentiels} (informations et processus
métier) et les \emph{actifs supports} (matériels, logiciels, réseaux,
personnes, locaux, prestataires) sur lesquels ils reposent.
\end{boxdef}

La distinction entre actifs essentiels et actifs supports est la clé de
toute démarche de sécurité~: on protège un serveur \emph{parce qu'}il porte
la facturation, et c'est la valeur de la facturation qui fixe l'effort.
L'ISO/IEC~27000 en donne le vocabulaire commun~\cite{iso27000_2018}.

\begin{tableaucours}{L{3.4cm}L{4.2cm}L{4.6cm}}{%
    \tableauHeader{Actif essentiel} &
    \tableauHeader{Actifs supports} &
    \tableauHeader{Touché le 5 mars~?}
}
Facturation & Application intégrée, serveurs de Douala et de Garoua &
Oui~: serveurs chiffrés \\
Données des clients & Base de données, partage de fichiers & Oui~: copie
probable de 40~Go \\
Suivi des expéditions & Portail web, base de données & Indirectement~:
données non mises à jour \\
Messagerie & Service en nuage & Non, d'après les premiers éléments \\
Sauvegardes & Partage réseau du domaine & Oui~: supprimées \\
\end{tableaucours}
\vspace{-0.8em}
\begin{center}{\small\itshape Tableau~1.1. Actifs de LiZbiZ et Opération GAROUA}\end{center}

% ============================================================================
\section{Sécurité de l'Information et Cybersécurité}
\label{sec:c01-securite}
% ============================================================================

La \termecle{sécurité de l'information} protège l'information sous toutes
ses formes~: numérique, papier, orale. La \termecle{sécurité informatique}
protège les systèmes qui la traitent. La \termecle{cybersécurité} désigne
la sécurité de ce qui passe par les réseaux numériques.

\begin{boxjuris}
\textbf{\loicam{n°~2010/012}, article~4, point~33} (texte officiel)~:
\textit{« Cybersécurité~: ensemble de mesures de prévention, de protection
et de dissuasion d'ordre technique, organisationnel, juridique, financier,
humain, procédural et autres actions permettant d'atteindre les objectifs
de sécurité fixés à travers les réseaux de communications électroniques,
les systèmes d'information et pour la protection de la vie privée des
personnes. »}

\medskip
\textbf{Analyse.} La définition légale insiste sur la diversité des
mesures~: la technique n'en est qu'une parmi sept. À Garoua, la faille
initiale était humaine (un mot de passe donné à une fausse page) et
contractuelle (aucune exigence imposée au prestataire), avant d'être
technique.
\end{boxjuris}

% ============================================================================
\section{Les Propriétés de Sécurité}
\label{sec:c01-proprietes}
% ============================================================================

\subsection{La triade}

Trois propriétés forment le socle, que la loi définit~:

\begin{boxjuris}
\textbf{\loicam{n°~2010/012}, article~4} (texte officiel)

\medskip
\textbf{Point~24}~: \textit{« Confidentialité~: maintien du secret des
informations et des transactions afin de prévenir la divulgation non
autorisée d'informations aux non destinataires permettant la lecture,
l'écoute, la copie illicite d'origine intentionnelle ou accidentelle durant
leur stockage, traitement ou transfert. »}

\medskip
\textbf{Point~48}~: \textit{« Intégrité des données~: critère de sécurité
définissant l'état d'un réseau de communications électroniques, d'un
système d'information ou d'un équipement terminal qui est demeuré intact et
permet de s'assurer que les ressources n'ont pas été altérées (modifiées ou
détruites) d'une façon tant intentionnelle qu'accidentelle, de manière à
assurer leur exactitude, leur fiabilité et leur pérennité. »}

\medskip
\textbf{Point~38}~: \textit{« Disponibilité~: critère de sécurité
permettant que les ressources des réseaux de communications électroniques,
des systèmes d'information ou des équipements terminaux soient accessibles
et utilisables selon les besoins (le facteur temps). »}
\end{boxjuris}

\subsection{Au-delà de la triade}

La triade ne suffit pas à décrire toutes les pertes. On lui ajoute
habituellement trois propriétés, et D.~Parker a proposé de remplacer la
triade par six attributs~\cite{parker1998}.

\begin{tableaucours}{L{3.0cm}L{5.0cm}L{4.2cm}}{%
    \tableauHeader{Propriété} &
    \tableauHeader{Question posée} &
    \tableauHeader{Mesures typiques}
}
Confidentialité & Seules les personnes autorisées peuvent-elles lire~? &
Contrôle d'accès, chiffrement \\
Intégrité & L'information est-elle exacte et complète~? & Hachage,
signatures, contrôles de saisie \\
Disponibilité & Est-elle accessible quand on en a besoin~? & Redondance,
sauvegardes, plan de continuité \\
Authenticité & L'origine est-elle bien celle que l'on croit~? &
Authentification, signature \\
Non-répudiation & L'auteur peut-il nier son action~? & Signature, journaux
protégés, horodatage \\
Traçabilité (imputabilité) & Peut-on savoir qui a fait quoi~? & Comptes
nominatifs, journalisation \\
\end{tableaucours}
\vspace{-0.8em}
\begin{center}{\small\itshape Tableau~1.2. Les propriétés de sécurité}\end{center}

\begin{boxjuris}
\textbf{\loicam{n°~2010/012}, article~4, point~58} (texte officiel)~:
\textit{« Non répudiation~: critère de sécurité assurant la disponibilité
de preuves qui peuvent être opposées à un tiers et utilisées pour prouver
la traçabilité d'une communication électronique qui a eu lieu. »}
\end{boxjuris}

\begin{boxscenario}{Opération GAROUA. Trois pertes, trois propriétés}
\textbf{Disponibilité}~: la facturation est arrêtée depuis le 5~mars à
21~h~40. \textbf{Confidentialité}~: 40~Go de dossiers clients ont
probablement été copiés. \textbf{Intégrité}~: tant que l'on ignore ce que
l'attaquant a modifié pendant trois jours de présence, on ne peut plus se
fier aux données de stock. \textbf{Traçabilité}~: le compte du technicien,
partagé entre plusieurs personnes chez le prestataire, ne permet pas de
savoir qui s'est réellement connecté.
\end{boxscenario}

% ============================================================================
\section{Exprimer des Besoins de Sécurité}
\label{sec:c01-besoins}
% ============================================================================

Un besoin de sécurité dit jusqu'où une propriété doit être garantie pour
un actif essentiel. On l'exprime sur une échelle commune, fixée par la
direction, souvent de~1 (faible) à~4 (critique)~: c'est la cotation
\textbf{DICP} (disponibilité, intégrité, confidentialité, preuve).

\begin{tableaucours}{L{3.6cm}cccc}{%
    \tableauHeader{Actif essentiel} &
    \tableauHeader{D} &
    \tableauHeader{I} &
    \tableauHeader{C} &
    \tableauHeader{P}
}
Facturation & 4 & 4 & 2 & 3 \\
Données des clients & 2 & 3 & 4 & 2 \\
Suivi des expéditions & 3 & 3 & 2 & 1 \\
Paiements des fournisseurs & 3 & 4 & 3 & 4 \\
Messagerie & 3 & 2 & 3 & 2 \\
\end{tableaucours}
\vspace{-0.8em}
\begin{center}{\small\itshape Tableau~1.3. Besoins de sécurité de LiZbiZ
(cotation proposée, à valider par la direction)}\end{center}

\begin{boxalert}
\textbf{Violation typique.} Coter tout à 4 « pour être tranquille ».

\smallskip
\textbf{Règle~:} une cotation sert à hiérarchiser. Si tout est critique,
rien ne l'est, et le budget se disperse. Chaque cote se justifie par
l'impact d'une perte~: combien de jours sans facturer l'entreprise
peut-elle supporter~?
\end{boxalert}

% ============================================================================
\section{Menace, Vulnérabilité, Risque}
\label{sec:c01-risque}
% ============================================================================

\begin{boxdef}
\termecle{Menace}~: cause potentielle d'un incident indésirable
(un groupe criminel, un employé négligent, une inondation).
\termecle{Vulnérabilité}~: faiblesse qu'une menace peut exploiter. La
\loicam{n°~2010/012} (art.~4, point~73) la définit comme un « défaut de
sécurité se traduisant soit intentionnellement, soit accidentellement par
une violation de la politique de sécurité ». \termecle{Risque}~: effet de
l'incertitude sur les objectifs~; en sécurité, la combinaison de la
vraisemblance d'un scénario et de son impact.
\end{boxdef}

\begin{tableaucours}{L{3.0cm}L{4.4cm}L{4.8cm}}{%
    \tableauHeader{Notion} &
    \tableauHeader{Opération GAROUA} &
    \tableauHeader{Ce que l'on peut y faire}
}
Menace & Groupe criminel qui vise une rançon & Peu de chose~: on ne
choisit pas ses attaquants \\
Vulnérabilités & Accès distant à un facteur~; correctif manquant~;
sauvegardes accessibles & Les corriger \\
Impact & Arrêt de la facturation, fuite de données & Le réduire
(sauvegardes hors ligne, plan de continuité) \\
Risque & Vraisemblance élevée, impact critique & Le traiter~: réduire,
transférer, éviter ou accepter (Chapitre~\ref{chap:risques}) \\
\end{tableaucours}
\vspace{-0.8em}
\begin{center}{\small\itshape Tableau~1.4. Les notions appliquées à l'Opération
GAROUA}\end{center}

On écrit souvent \(\text{risque} = \text{vraisemblance} \times
\text{impact}\). C'est une convention de cotation, utile pour classer,
dont le Chapitre~\ref{chap:risques} montre les limites.

% ============================================================================
\section{Les Principes de Conception}
\label{sec:c01-principes}
% ============================================================================

En 1975, J.~Saltzer et M.~Schroeder ont énoncé huit principes de
conception des mécanismes de protection, qui restent la référence de tout
architecte~\cite{saltzer1975}.

\begin{tableaucours}{L{4.0cm}L{8.2cm}}{%
    \tableauHeader{Principe} &
    \tableauHeader{Ce qu'il demande}
}
Économie de mécanisme & Faire simple et petit~: ce qui est simple se
vérifie \\
Valeurs par défaut sûres & Refuser par défaut~; n'autoriser
qu'explicitement \\
Médiation complète & Vérifier chaque accès, pas seulement le premier \\
Conception ouverte & La sécurité ne repose pas sur le secret du
mécanisme, mais sur celui des clés \\
Séparation des privilèges & Exiger deux conditions plutôt qu'une (par
exemple deux facteurs) \\
Moindre privilège & Ne donner que les droits nécessaires à la tâche \\
Mécanisme le moins partagé & Limiter ce qui est commun à plusieurs
utilisateurs \\
Acceptabilité psychologique & Une mesure trop pénible sera contournée \\
\end{tableaucours}
\vspace{-0.8em}
\begin{center}{\small\itshape Tableau~1.5. Les huit principes de Saltzer et
Schroeder}\end{center}

\begin{boxobjet}{M-01 · Double authentification de l'accès distant}
\textbf{Objectif~:} empêcher qu'un mot de passe volé suffise à entrer.
\textbf{Menace couverte~:} hameçonnage, réutilisation de mots de passe.
\textbf{Mise en œuvre~:} un second facteur (application, clé physique)
exigé pour toute connexion distante, prestataires compris.
\textbf{Principe~:} séparation des privilèges.
\textbf{Vérification~:} tentative de connexion avec le seul mot de passe,
qui doit échouer~; revue mensuelle des comptes distants.
\end{boxobjet}

\begin{boxmetier}{RSSI}
\textbf{Poser les mots justes.} Le RSSI traduit un incident en propriétés
atteintes et en actifs touchés~: « la disponibilité de la facturation est
perdue depuis 34 heures~; la confidentialité des données clients est
probablement atteinte~; l'intégrité des stocks est inconnue ». C'est la
seule forme qui permette à une direction de décider.
\end{boxmetier}

\begin{boxmetier}{Développeur}
\textbf{Appliquer les principes dès la conception.} Valeurs par défaut
sûres (un nouveau compte n'a aucun droit), médiation complète (chaque
requête de l'application vérifie les droits), conception ouverte (aucun
mot de passe dans le code).
\end{boxmetier}

\separateur

\begin{boxresume}
\textbf{Ce qu'il faut retenir de ce chapitre~:}
\begin{itemize}
    \item \textbf{On protège des actifs essentiels} (informations,
          processus) à travers leurs actifs supports.
    \item \textbf{La cybersécurité} combine des mesures techniques,
          organisationnelles, juridiques et humaines (loi n°~2010/012,
          art.~4).
    \item \textbf{Six propriétés}~: confidentialité, intégrité,
          disponibilité, authenticité, non-répudiation, traçabilité.
    \item \textbf{Les besoins} se cotent sur une échelle commune, pour
          hiérarchiser.
    \item \textbf{Menace, vulnérabilité, impact} composent le risque~; on
          agit sur les vulnérabilités et l'impact.
    \item \textbf{Huit principes de conception} guident toute mesure.
\end{itemize}
\end{boxresume}

\begin{boxinfo}
\textbf{Points de contrôle de l'élève.} Je sais~: distinguer actif
essentiel et actif support~; citer et définir les six propriétés~; coter
un actif en DICP~; distinguer menace et vulnérabilité~; nommer les huit
principes de Saltzer et Schroeder et en illustrer trois.
\end{boxinfo}

\bigskip

\begin{boxexo}[title={\ding{45}\quad Exercice~1.1. Essentiel ou support~? \niveaubase}]
Classez~: le fichier des clients~; le serveur de base de données~; le
processus de facturation~; le technicien d'astreinte~; le contrat avec
SAHEL INFOGÉRANCE~; la liaison entre Douala et Garoua.
\end{boxexo}

\begin{boxexo}[title={\ding{45}\quad Exercice~1.2. Quelle propriété~? \niveaubase}]
Dites quelle propriété est atteinte~: (a) un virement de 5~millions est
modifié en 50~millions~; (b) un fournisseur nie avoir envoyé une facture~;
(c) un concurrent lit les tarifs de LiZbiZ~; (d) le portail de suivi est
saturé de requêtes~; (e) une action est faite avec un compte partagé.
\end{boxexo}

\begin{boxexo}[title={\ding{45}\quad Exercice~1.3. Le texte de la loi \niveaubase}]
Relevez, dans la définition légale de la confidentialité, les trois
moments où l'information doit être protégée, et donnez pour chacun une
mesure.
\end{boxexo}

\begin{boxexo}[title={\ding{45}\quad Exercice~1.4. Coter \niveaumoyen}]
Cotez en DICP trois actifs essentiels d'une agence bancaire de la BANQUE
SANAGA, et justifiez chaque cote de~4.
\end{boxexo}

\begin{boxexo}[title={\ding{45}\quad Exercice~1.5. Menace ou vulnérabilité~? \niveaumoyen}]
Classez~: un mot de passe par défaut~; une coupure d'électricité~; un
employé mécontent~; un serveur sans correctif~; une crue du Wouri~; une
porte de salle serveur laissée ouverte.
\end{boxexo}

\begin{boxexo}[title={\ding{45}\quad Exercice~1.6. Les principes à Garoua \niveaumoyen}]
Pour trois défaillances de l'Opération GAROUA, indiquez le principe de
Saltzer et Schroeder qui n'a pas été respecté.
\end{boxexo}

\begin{boxexo}[title={\ding{45}\quad Exercice~1.7. Une mesure \niveaumoyen}]
Rédigez, sur le modèle de la mesure M-01, la fiche d'une mesure qui
protège les sauvegardes contre un rançongiciel.
\end{boxexo}

\begin{boxexo}[title={\ding{45}\quad Exercice~1.8. Conception ouverte \niveaumoyen}]
Un développeur de LiZbiZ propose de « cacher » l'adresse de la page
d'administration au lieu de la protéger par un mot de passe. Répondez-lui
avec le principe de conception ouverte.
\end{boxexo}

\begin{boxexo}[title={\ding{45}\quad Exercice~1.9. Acceptabilité \niveauexpert}]
LiZbiZ impose un mot de passe de 20~caractères changé chaque semaine. Six
mois plus tard, l'auditeur trouve des mots de passe écrits sur des
post-it. Analysez la situation avec les principes du chapitre et proposez
une mesure qui protège mieux.
\end{boxexo}

\begin{boxexo}[title={\ding{45}\quad Exercice~1.10. La note du lundi \niveauexpert}]
Réécrivez la phrase de M.~ETEME, « la sécurité a été percée », en une
note de cinq lignes qui nomme les actifs touchés, les propriétés atteintes
et ce qui reste inconnu.
\end{boxexo}

\bigskip

\begin{boxinfo}
\textbf{Pour aller plus loin.}
\begin{itemize}
    \item ISO/IEC~27000, vocabulaire de la sécurité de
          l'information~\cite{iso27000_2018}.
    \item J.~Saltzer et M.~Schroeder, \emph{The Protection of Information
          in Computer Systems}~\cite{saltzer1975}.
    \item R.~Anderson, \emph{Security Engineering}~\cite{anderson2020}.
\end{itemize}
\end{boxinfo}

\begin{boxinfo}
\textbf{Transition.} On sait maintenant ce que l'on protège. Le
Chapitre~\ref{chap:menaces} décrit contre quoi~: les attaquants, leurs
techniques et la chaîne qu'ils suivent.
\end{boxinfo}
